% =====================================================================
% APENDICES -- Roteiros padronizados do protocolo experimental.
%
% Destinados a instruir a submissao ao CEP e a garantir uniformidade entre
% sessoes e entre participantes. Convencoes de leitura:
%
%   [A] ...  -> ACAO do pesquisador (nao e falado)
%   Fala em bloco recuado e italico -> texto LIDO EM VOZ ALTA, sem improviso
%   [D] ...  -> ponto de DECISAO, com criterio explicito
%
% Alinhados ao desenho de 4 condicoes (C1..C4) do Capitulo 3.
% =====================================================================

\chapter{Roteiro Geral da Sessão Experimental}\label{ap:roteiro-geral}

Este apêndice reúne o roteiro comum às quatro sessões experimentais, aplicável independentemente da condição designada. As falas em itálico são lidas em voz alta, sem improviso, para que todos os participantes recebam instrução idêntica. As ações do pesquisador aparecem entre colchetes e não são verbalizadas.

\section*{Etapa 1 --- Recepção e verificação de conformidade (5 min)}

\noindent\textbf{[A]} Receber o participante, confirmar sua identificação por código e conferir se a sessão corresponde à sequência designada.

\begin{quote}\itshape
Olá, tudo bem? Obrigado por comparecer. Esta é a sua sessão número [\underline{\hspace{1cm}}] de quatro. Antes de começar, preciso fazer algumas perguntas rápidas sobre as últimas horas, para garantir que as medidas fiquem comparáveis entre as sessões. Não existe resposta certa ou errada, e nada aqui afeta a sua participação.
\end{quote}

\noindent\textbf{[A]} Aplicar a lista de conformidade pré-sessão: horas de sono, qualidade do sono, consumo de cafeína, nicotina e álcool nas últimas 24 horas, exercício físico recente, doença ou dor atual, mudança de medicação e eventos estressantes antes de chegar ao laboratório.

\noindent\textbf{[D]} Critério de reagendamento: consumo de cafeína ou nicotina nas últimas 2 horas, consumo de álcool nas últimas 24 horas, exercício vigoroso nas últimas 12 horas ou doença aguda em curso. Havendo qualquer desses fatores, a sessão é reagendada e o motivo é registrado.

\begin{quote}\itshape
Vou também perguntar se você sentiu algum desconforto depois da sessão anterior, como dor de cabeça, enjoo, cansaço nos olhos ou dor muscular.
\end{quote}

\noindent\textbf{[A]} Registrar eventos adversos relatados. Encaminhar ao roteiro de eventos adversos do Apêndice~\ref{ap:encerramento} se houver relato de sintoma moderado ou grave.

\section*{Etapa 2 --- Colocação dos sensores (5 min)}

\begin{quote}\itshape
Agora vou colocar os sensores. São dois. O primeiro é uma cinta elástica que fica em volta do tórax e mede os batimentos do coração; você mesmo vai colocá-la, em privacidade, e eu explico como. O segundo são dois eletrodos pequenos que ficam nos dedos da sua mão [\underline{esquerda/direita}], a que você não usa para jogar, e medem a atividade da pele. Nenhum dos dois emite corrente elétrica; os dois apenas registram. Se algo incomodar, me avise a qualquer momento.
\end{quote}

\noindent\textbf{[A]} Orientar a colocação da cinta peitoral em privacidade. Umedecer os eletrodos conforme especificação. Fixar os eletrodos de \gls{eda} nas falanges intermediárias dos dedos indicador e médio da mão não dominante.

\noindent\textbf{[A]} Verificar a qualidade dos sinais por no mínimo 60 segundos antes de prosseguir. Registrar temperatura e umidade da sala.

\noindent\textbf{[D]} Critério de reposicionamento: sinal cardíaco com perda de batimentos ou sinal de \gls{eda} fora da faixa esperada. Persistindo após duas tentativas, registrar como falha de sensor e decidir sobre a continuidade da sessão.

\section*{Etapa 3 --- Linha de base em repouso (10 min)}

\begin{quote}\itshape
Agora vem um período de dez minutos de repouso. Você vai ficar sentado, confortável, com os pés apoiados no chão e as mãos apoiadas. Peço que evite conversar, mexer no celular e se movimentar mais do que o necessário. Pode manter os olhos abertos, olhando para um ponto à sua frente. Eu vou ficar em silêncio e aviso quando terminar.
\end{quote}

\noindent\textbf{[A]} Marcar o início da linha de base. Permanecer em silêncio, sem interagir. Os cinco minutos finais, livres de artefato, constituem a janela de análise primária.

\section*{Etapa 4 --- Instrumentos pré-estressor (3 a 5 min)}

\begin{quote}\itshape
Vou pedir agora que você responda a algumas escalas curtas sobre como se sente \emph{neste momento}. Não pense muito em cada uma; a primeira impressão é o que interessa. Responda com sinceridade, lembrando que suas respostas são anônimas.
\end{quote}

\noindent\textbf{[A]} Aplicar, nesta ordem: escalas visuais analógicas de ansiedade, estresse e calma; \gls{sam}; subescala de tensão do \gls{brums}. Na primeira sessão, aplicar também o \gls{ssq} de base.

% =====================================================================
\chapter{Roteiro da Tarefa de Indução de Estresse}\label{ap:estressor}

Este apêndice descreve a tarefa de aritmética mental sob pressão de tempo, aplicada antes de cada condição experimental, com duração total de cinco minutos. O procedimento é o mesmo nas quatro sessões, variando apenas os parâmetros numéricos.

\section*{Fundamento e cuidado ético}

A tarefa é deliberadamente difícil e foi calibrada para que o desempenho fique abaixo do confortável. Essa característica é o mecanismo da indução e, por isso, é comunicada ao participante no \gls{tcle} e novamente no \textit{debriefing} ao fim de cada sessão. Não há engano quanto à natureza da tarefa: o participante sabe, desde o consentimento, que enfrentará uma atividade difícil sob tempo. A retroalimentação é impessoal, referida apenas ao acerto ou erro do item, sem comparação com outros participantes, sem julgamento de capacidade e sem qualquer sugestão de consequência.

\section*{Instrução (2 min, incluindo bloco de prática)}

\begin{quote}\itshape
Agora você vai fazer uma tarefa de cálculo mental no computador. A tarefa mostra uma conta de subtração e você digita o resultado. Assim que responder, aparece a próxima. Há um tempo limite para cada conta, mostrado por uma barra na tela, e a tela indica se a resposta foi certa ou errada.

Duas coisas importantes. Primeira: a tarefa é difícil de propósito e o tempo é curto de propósito. É esperado que você erre várias contas; isso faz parte do desenho e não é uma avaliação da sua capacidade. Segunda: o que interessa para a pesquisa é a sua reação a esse tipo de situação, e não a sua pontuação. Ninguém verá seu desempenho individual.

Vamos fazer um bloco curto de prática para você se familiarizar. Alguma dúvida?
\end{quote}

\noindent\textbf{[A]} Executar o bloco de prática (aproximadamente 30 segundos, sem registro). Esclarecer dúvidas operacionais. Não antecipar a dificuldade do bloco principal além do já dito.

\section*{Execução (5 min)}

\begin{quote}\itshape
Agora começa a tarefa que vale para a pesquisa. Ela dura cinco minutos. Procure acertar o máximo que conseguir dentro do tempo de cada conta. Vou permanecer na sala, em silêncio.
\end{quote}

\noindent\textbf{[A]} Marcar o início da tarefa. Permanecer na sala, sem interagir e sem observar ostensivamente o participante.

\noindent\textbf{[A]} Parâmetros da tarefa: subtrações seriadas encadeadas, com limite de tempo por item ajustado para manter a taxa de acerto abaixo do desempenho confortável. Os valores iniciais e o passo variam entre as quatro sessões, preservando a dificuldade. \todo{fixar os parametros numericos exatos por sessao apos o estudo-piloto e registra-los aqui, para que o roteiro fique completo na submissao ao CEP.}

\noindent\textbf{[D]} Critério de interrupção: manifestação de desconforto acentuado, choro, tremor, solicitação de parada ou relato de mal-estar. Havendo qualquer um, interromper imediatamente, aplicar o roteiro de acolhimento do Apêndice~\ref{ap:encerramento} e registrar como evento adverso.

\section*{Verificação de manipulação (2 a 3 min)}

\begin{quote}\itshape
Terminou. Antes de seguir, responda de novo às escalas sobre como você se sente \emph{agora}.
\end{quote}

\noindent\textbf{[A]} Reaplicar as escalas visuais analógicas e o \gls{sam}. Esses valores constituem a medida pós-estressor, da qual se mede a recuperação.

% =====================================================================
\chapter{Roteiros das Condições Experimentais}\label{ap:condicoes}

Cada participante realiza uma condição por sessão, na ordem definida pelo quadrado de Williams. A exposição dura 15 minutos nas quatro condições. As instruções abaixo são lidas imediatamente antes da exposição.

\section{Condição C1 --- \textit{Cozy Pong} Relaxante}

\noindent\textbf{[A]} Higienizar o \gls{hmd} e a interface facial. Ajustar o dispositivo, verificar nitidez e conferir o volume previamente calibrado na sessão de orientação.

\begin{quote}\itshape
Nesta sessão você vai jogar o \textit{Cozy Pong}. É um jogo de realidade virtual em que bolas se aproximam no ritmo da música e você as rebate com a raquete que está na sua mão. A música é calma e as bolas vêm em um ritmo tranquilo e previsível.

Três observações. Você não pode perder: não existe game over nem punição por errar. Não é preciso se apressar nem fazer movimentos amplos; um movimento suave é suficiente. E não há meta de pontuação, então não se preocupe com o placar.

A ideia é simplesmente jogar de forma tranquila por quinze minutos. Se sentir qualquer desconforto, enjoo ou tontura, avise imediatamente e paramos na hora. Podemos começar?
\end{quote}

\noindent\textbf{[A]} Marcar o início da exposição. Permanecer próximo, em silêncio, atento a sinais de desconforto e à segurança do espaço.

\section{Condição C2 --- \textit{Cozy Pong} Animado}

\noindent\textbf{[A]} Procedimento de higienização e ajuste idêntico ao de C1.

\begin{quote}\itshape
Nesta sessão você vai jogar o \textit{Cozy Pong} em uma versão mais movimentada. A música é mais acelerada, as bolas vêm com mais frequência e mais rápido, e há um limite de erros: se você errar muitas vezes, a partida termina antes do tempo.

O objetivo é acompanhar o ritmo e acertar o máximo que conseguir. Se a partida terminar antes dos quinze minutos, ela recomeça automaticamente até completar o tempo.

Se sentir qualquer desconforto, enjoo ou tontura, avise imediatamente e paramos na hora. Podemos começar?
\end{quote}

\noindent\textbf{[A]} Marcar o início da exposição. Registrar eventuais reinícios de partida na telemetria.

\section{Condição C3 --- Audição da Trilha}

\noindent\textbf{[A]} Acomodar o participante sentado, na mesma poltrona e postura da linha de base. Colocar os fones e conferir o volume, idêntico ao empregado em C1.

\begin{quote}\itshape
Nesta sessão não há jogo nem óculos. Você vai apenas ouvir música por quinze minutos, sentado, com os fones. É a mesma trilha que aparece em uma das outras sessões.

Não há nenhuma tarefa: você não precisa prestar atenção em nada específico, não vou fazer perguntas sobre a música depois e não há nada a memorizar. Pode manter os olhos abertos ou fechados, como preferir. Peço apenas que evite mexer no celular, conversar ou dormir.

Eu vou ficar em silêncio e aviso quando terminar. Podemos começar?
\end{quote}

\noindent\textbf{[A]} Marcar o início da exposição. Registrar manualmente início e fim, conforme procedimento padronizado.

\noindent\textbf{[D]} Se o participante adormecer, despertá-lo com delicadeza, registrar o episódio e anotar o intervalo aproximado, que será considerado na análise.

\section{Condição C4 --- Meditação Guiada em Realidade Virtual}

\noindent\textbf{[A]} Procedimento de higienização e ajuste idêntico ao de C1. Carregar a cena de meditação do mesmo aplicativo, com a trilha em volume reduzido padronizado.

\begin{quote}\itshape
Nesta sessão você vai usar os óculos, mas não há jogo. Você verá um ambiente virtual tranquilo e ouvirá uma narração que guia um exercício de respiração, por quinze minutos.

Você fica sentado, sem precisar se mover nem usar os controles. Basta seguir as instruções da narração no seu próprio ritmo. Não existe forma certa ou errada de fazer, e é completamente normal a atenção se dispersar; quando perceber que isso aconteceu, simplesmente volte à respiração.

Se sentir qualquer desconforto, enjoo ou tontura, avise imediatamente e paramos na hora. Podemos começar?
\end{quote}

\noindent\textbf{[A]} Marcar o início da exposição. Permanecer em silêncio durante toda a sessão.

\subsection*{Roteiro da narração guiada}

O texto abaixo é a narração da condição C4, escrita especificamente para este protocolo e, portanto, livre de restrição de licenciamento de terceiros. É gravada previamente, em áudio único, com voz calma e andamento regular, e reproduzida de forma idêntica para todos os participantes. Os intervalos de silêncio entre os blocos são parte do roteiro e devem ser respeitados na gravação.

\begin{description}[style=nextline,leftmargin=1.6cm]

\item[\textbf{Bloco 1 (0:00--2:00) --- Acomodação.}]
\textit{Encontre uma posição confortável na cadeira. Deixe os pés apoiados no chão e as mãos descansando sobre as pernas. Se for confortável para você, feche os olhos; se preferir mantê-los abertos, apenas deixe o olhar suave, sem focar em nada específico. \\[4pt] Não há nada que você precise fazer nos próximos minutos. Nenhuma tarefa, nenhum objetivo. Apenas um tempo para observar. \\[4pt] Comece percebendo os pontos em que seu corpo toca a cadeira. As costas apoiadas. Os pés no chão. O peso dos braços. \textnormal{[silêncio de 20 segundos]}}

\item[\textbf{Bloco 2 (2:00--5:00) --- Atenção à respiração.}]
\textit{Agora leve a atenção para a sua respiração. Não é preciso mudar nada nela. Apenas perceba o ar entrando e o ar saindo, no ritmo em que já está acontecendo. \\[4pt] Talvez você note o ar mais fresco ao entrar pelas narinas, e mais morno ao sair. Talvez perceba o peito ou o abdômen subindo e descendo. Escolha o ponto onde a respiração é mais fácil de notar e deixe a atenção repousar ali. \textnormal{[silêncio de 30 segundos]} \\[4pt] Se em algum momento você perceber que a atenção foi para outro lugar, para um pensamento, uma lembrança, um som da sala, está tudo bem. Isso vai acontecer várias vezes, e não é um erro. Quando notar, apenas volte, com gentileza, para a respiração. \textnormal{[silêncio de 40 segundos]}}

\item[\textbf{Bloco 3 (5:00--9:00) --- Respiração com contagem.}]
\textit{Vamos agora dar um ritmo um pouco mais lento à respiração. Vou contar junto com você nas primeiras vezes. \\[4pt] Inspire contando até quatro: um, dois, três, quatro. Agora solte o ar, mais devagar, contando até seis: um, dois, três, quatro, cinco, seis. \\[4pt] De novo. Inspire: um, dois, três, quatro. Expire: um, dois, três, quatro, cinco, seis. \\[4pt] Continue nesse ritmo por conta própria. Se o ritmo for desconfortável, ajuste para o que for natural para você. O importante é que a saída do ar seja um pouco mais longa que a entrada. \textnormal{[silêncio de 60 segundos]} \\[4pt] Siga assim, no seu tempo. \textnormal{[silêncio de 90 segundos]}}

\item[\textbf{Bloco 4 (9:00--12:30) --- Varredura corporal breve.}]
\textit{Deixe a respiração voltar ao ritmo natural, sem contagem. \\[4pt] Agora leve a atenção para os ombros. Perceba se há alguma tensão ali e, se houver, veja se ela pode ceder um pouco, junto com a próxima expiração. \textnormal{[silêncio de 20 segundos]} \\[4pt] Leve a atenção para a mandíbula e para a região ao redor dos olhos. São lugares onde a tensão costuma se acumular sem que a gente perceba. \textnormal{[silêncio de 20 segundos]} \\[4pt] Agora para as mãos. Perceba a temperatura, o contato com as pernas, qualquer sensação que esteja presente. \textnormal{[silêncio de 20 segundos]} \\[4pt] E agora o corpo inteiro, de uma vez, sentado aqui, respirando. \textnormal{[silêncio de 40 segundos]}}

\item[\textbf{Bloco 5 (12:30--15:00) --- Encerramento.}]
\textit{Vamos começar a encerrar. \\[4pt] Volte a atenção para a respiração por mais alguns instantes. \textnormal{[silêncio de 30 segundos]} \\[4pt] Agora perceba novamente os pontos de apoio do corpo na cadeira, os pés no chão. \\[4pt] Perceba os sons ao redor. \\[4pt] Quando estiver pronto, faça um movimento leve com os dedos das mãos, e depois com os ombros. \\[4pt] Se seus olhos estiverem fechados, abra-os devagar. A sessão terminou. Obrigado.}

\end{description}

\noindent\textbf{[A]} A cena visual que acompanha a narração é sóbria e estática, sem elementos de jogo, sem pontuação e sem movimento súbito de câmera, de modo a não introduzir estimulação visual que compita com a instrução. \todo{produzir a cena no projeto Unity e gravar o audio da narracao acima; registrar aqui a duracao final exata de cada bloco.}

% =====================================================================
\chapter{Roteiro de Encerramento, \textit{Debriefing} e Segurança}\label{ap:encerramento}

\section*{Etapa 9 --- Instrumentos pós-intervenção (5 a 8 min)}

\begin{quote}\itshape
Terminamos a atividade. Agora vou pedir que você responda a mais um conjunto de escalas, novamente sobre como se sente \emph{neste momento}, e algumas perguntas sobre como foi a experiência.
\end{quote}

\noindent\textbf{[A]} Aplicar: escalas visuais analógicas, \gls{sam}, \gls{brums}, \gls{tlx}, Borg CR10 e itens de fruição. Em C1, C2 e C4, acrescentar \gls{ipq} e \gls{ssq}. Em C1 e C2, acrescentar os itens de sincronização percebida e as questões abertas.

\section*{Etapa 10 --- Recuperação (10 min)}

\begin{quote}\itshape
Agora vem um período de dez minutos de repouso, igual ao do começo da sessão. Sentado, confortável, sem conversar e sem mexer no celular. Eu fico em silêncio e aviso quando terminar.
\end{quote}

\noindent\textbf{[A]} Marcar o início da recuperação. Os cinco minutos finais, livres de artefato, constituem a janela de análise primária de \gls{vfc} de recuperação.

\section*{Etapa 11 --- Medidas finais (2 a 3 min)}

\noindent\textbf{[A]} Reaplicar as escalas visuais analógicas e o \gls{sam}. Aplicar a verificação de eventos adversos.

\section*{Etapa 12 --- \textit{Debriefing} e liberação (3 a 5 min)}

O \textit{debriefing} é obrigatório ao fim de \emph{todas} as sessões, inclusive quando o participante não manifesta desconforto.

\begin{quote}\itshape
Antes de você ir, quero explicar uma coisa sobre a tarefa de cálculo do início da sessão.

Aquela tarefa foi montada de propósito para ser difícil e para ter um tempo curto. A dificuldade foi calibrada para que praticamente todo mundo erre uma boa parte das contas, independentemente de habilidade com matemática. Ou seja: se você errou bastante, isso era o esperado e não diz nada sobre você. O objetivo dela era gerar, por alguns minutos, aquela sensação de pressão que a gente sente no dia a dia, para que eu pudesse medir como diferentes atividades ajudam a pessoa a se recuperar depois.

Seu desempenho na tarefa não é analisado individualmente e não será mostrado a ninguém.

Como você está se sentindo agora?
\end{quote}

\noindent\textbf{[A]} Aguardar a resposta e acolher. Retirar os equipamentos. Confirmar que o participante está bem antes da saída.

\noindent\textbf{[D]} Critério de retenção: se o participante relatar desconforto persistente, permanecer com ele, oferecer água e ambiente tranquilo, e liberar somente após a normalização do relato. Registrar como evento adverso e, se apropriado, entregar a informação de encaminhamento para apoio psicológico.

\begin{quote}\itshape
Sua próxima sessão está marcada para [\underline{\hspace{2cm}}]. Lembrando das orientações: evitar cafeína e nicotina nas duas horas anteriores, álcool nas 24 horas anteriores e exercício intenso nas 12 horas anteriores. Se precisar remarcar, é só avisar. Obrigado pela participação.
\end{quote}

\section*{Critérios de interrupção}

A sessão é interrompida imediatamente, em qualquer etapa, nas seguintes situações:

\begin{enumerate}
    \item solicitação do participante, por qualquer motivo e sem necessidade de justificativa;
    \item relato de náusea, tontura, desorientação ou desconforto visual moderado ou intenso durante a exposição à \gls{rv};
    \item manifestação de sofrimento acentuado durante a tarefa de indução de estresse;
    \item dor, desconforto físico ou fadiga excessiva;
    \item qualquer intercorrência de saúde.
\end{enumerate}

Em todos os casos, aplica-se o roteiro de acolhimento, registra-se o evento adverso e avalia-se, com o participante, a continuidade no estudo. A interrupção não implica exclusão automática dos dados já coletados nas demais sessões.

\section*{Procedimentos de higiene}

\noindent\textbf{[A]} Entre participantes: higienizar o \gls{hmd}, a interface facial e os controles com produto apropriado; substituir a proteção facial descartável; higienizar a cinta peitoral conforme especificação do fabricante; utilizar eletrodos descartáveis ou higienizá-los conforme protocolo. Registrar a execução da higienização em ficha própria.
